\documentclass{amsart}
\usepackage{amsmath,amssymb,amsthm,hyperref}
\theoremstyle{plain}
\newtheorem{theorem}{Theorem}
\newtheorem{lemma}{Lemma}
\newtheorem{proposition}{Proposition}
\theoremstyle{definition}
\newtheorem{definition}{Definition}
\newtheorem{remark}{Remark}
\begin{document}

\title{A coherent Bayesian resolution of the two-envelopes and necktie paradoxes}
\maketitle

\section{Formalization and statement of the resolution}

We work entirely inside subjective expected-utility theory. An agent's beliefs
about the relevant monetary quantities are encoded by a prior measure $\pi$ on
$(0,\infty)$ (or on $(0,\infty)^2$), and the agent acts so as to maximize
subjective expected utility. We assume throughout, as is standard in the
statements of both puzzles, that utility is linear in money (risk-neutrality);
this is to isolate the genuine source of the paradoxes, which is
\emph{probabilistic}, not a matter of nonlinearity of utility. (Remark~\ref{rem:utility}
records what changes under nonlinear utility.)

The resolution we will prove has three parts.

\begin{enumerate}
\item[(R1)] \textbf{The defective step is the conditional symmetry assumption.}
In both puzzles the naive computation silently assumes that, conditional on the
information actually available to the agent (the observed value $x$ in the
envelope problem; the observed value $v$ of one's own tie in the necktie
problem), the two states ``mine is the larger'' and ``mine is the smaller'' are
\emph{equiprobable, for every observed value}. We prove (Lemma~\ref{lem:noproper})
that \emph{no} proper prior has this property simultaneously for all observed
values. The naive reasoning therefore is not the evaluation of any coherent
posterior; it is the misapplication of the principle of indifference to a
\emph{conditional} probability where indifference does not hold.

\item[(R2)] \textbf{If one nonetheless forces conditional symmetry, the
underlying prior is improper and the relevant expectation is undefined.} The
unique belief structure making the two states equiprobable at every observed
value is the scale-invariant (log-uniform) improper prior. Under it the marginal
expected value of money is infinite (Lemma~\ref{lem:improper}), so the
``always switch'' conclusion rests on manipulating an undefined / divergent
expectation, where the law of total expectation fails. This is illegitimate
conditioning relative to an improper measure.

\item[(R3)] \textbf{The general normative principle.} For any proper prior the
gain from exchange has \emph{zero unconditional expectation} whenever the prior
is symmetric under interchange of the two parties
(Theorem~\ref{thm:fubini}). More generally, the only legitimate quantity on
which to base a decision is $\mathbb E_\pi[\,U(\text{after})-U(\text{before})\,]$,
computed from a prior under which this expectation is well defined (finite) and
the law of total expectation (Fubini--Tonelli) applies. Under this principle
the symmetric exchange is never strictly advantageous to both parties, and
``always switch'' is never optimal for all observed values.
\end{enumerate}

The single coherent normative principle covering all exchange-type problems is
thus:

\begin{quote}
\emph{A decision must maximize the expectation of the change in utility taken
under a (proper, or at least integrable) prior for which that expectation is
well defined; conditional probabilities used in the computation must be the
genuine posteriors $\pi(\,\cdot\mid \text{observed data})$, never values
assigned by an unexamined appeal to symmetry. Whenever the prior is symmetric
under the exchange of roles, the expected change in utility from exchanging is
zero, so the exchange cannot be advantageous to both parties.}
\end{quote}

\section{The two-envelopes problem}

\subsection{Setup} Let $T>0$ denote the smaller of the two amounts, so the two
envelopes contain $T$ and $2T$. The agent's beliefs about $T$ are a prior; let
its law have density $f$ on $(0,\infty)$ with respect to Lebesgue measure
(the discrete case is identical with sums replacing integrals). The agent picks
one envelope uniformly at random; let $X$ be its content and $Y$ the other
envelope's content. Thus the joint law of $(T,\text{side})$ is: $T\sim f$, and
independently the agent holds the small or the large envelope each with
probability $\tfrac12$. Consequently
\begin{equation}\label{eq:Xlaw}
\Pr(X=T)=\Pr(X=2T)=\tfrac12 .
\end{equation}

\subsection{The correct posterior} Fix an observed value $x>0$. Conditioning on
$X=x$, exactly two hypotheses are consistent with the observation:
\[
H_{\text{small}}:\ T=x\ \ (\text{so }Y=2x),
\qquad
H_{\text{large}}:\ T=x/2\ \ (\text{so }Y=x/2).
\]
The (unnormalized) likelihood-weighted prior masses of these hypotheses are,
respectively, $\tfrac12 f(x)$ (the pair is $(x,2x)$ and the agent drew the small
side) and $\tfrac12 f(x/2)$ (the pair is $(x/2,x)$ and the agent drew the large
side). By Bayes' rule the common factor $\tfrac12$ cancels and
\begin{equation}\label{eq:posterior}
\Pr(Y=2x\mid X=x)=\frac{f(x)}{f(x)+f(x/2)},
\qquad
\Pr(Y=x/2\mid X=x)=\frac{f(x/2)}{f(x)+f(x/2)},
\end{equation}
valid for every $x$ with $f(x)+f(x/2)>0$. Hence the correct conditional
expectation is
\begin{equation}\label{eq:condexp}
\mathbb E[Y\mid X=x]
=\frac{2x\,f(x)+\tfrac{x}{2}\,f(x/2)}{f(x)+f(x/2)} .
\end{equation}

\begin{lemma}[The naive value is generally false]\label{lem:naive}
The naive formula $\mathbb E[Y\mid X=x]=\tfrac54 x$ holds for a given $x$ if and
only if $f(x)=f(x/2)$. Switching is strictly favorable at $x$
\(\bigl(\mathbb E[Y\mid X=x]>x\bigr)\) if and only if $2f(x)>f(x/2)$, and is
strictly unfavorable if and only if $2f(x)<f(x/2)$.
\end{lemma}

\begin{proof}
From \eqref{eq:condexp},
$\mathbb E[Y\mid X=x]-\tfrac54 x
=\dfrac{(2x-\tfrac54x)f(x)+(\tfrac x2-\tfrac54x)f(x/2)}{f(x)+f(x/2)}
=\dfrac{\tfrac34 x\,f(x)-\tfrac34 x\,f(x/2)}{f(x)+f(x/2)}
=\dfrac{\tfrac34 x\,\bigl(f(x)-f(x/2)\bigr)}{f(x)+f(x/2)}$,
which vanishes iff $f(x)=f(x/2)$. Likewise
\[
\mathbb E[Y\mid X=x]-x
=\frac{(2x-x)f(x)+(\tfrac x2-x)f(x/2)}{f(x)+f(x/2)}
=\frac{x\,f(x)-\tfrac x2 f(x/2)}{f(x)+f(x/2)}
=\frac{\tfrac x2\,\bigl(2f(x)-f(x/2)\bigr)}{f(x)+f(x/2)} ,
\]
whose sign is that of $2f(x)-f(x/2)$ (as $x>0$ and the denominator is positive).
\end{proof}

The naive computation $\tfrac12(2x)+\tfrac12(x/2)$ tacitly sets
$\Pr(Y=2x\mid X=x)=\Pr(Y=x/2\mid X=x)=\tfrac12$, i.e. by \eqref{eq:posterior}
it assumes $f(x)=f(x/2)$. This is the defective step (R1): \eqref{eq:Xlaw} gives
the \emph{unconditional} side-symmetry $\Pr(X=T)=\Pr(X=2T)=\tfrac12$, but
side-symmetry of $X$ versus $T$ is \emph{not} the same as posterior symmetry of
$Y$ versus $x$ at a fixed observed $x$, because observing the value $x$ carries
information about whether $x$ is more plausibly the small or the large amount.

\subsection{No proper prior makes switching always favorable}

\begin{lemma}[Conditional symmetry forces an improper prior]\label{lem:noproper}
There is no probability density $f$ on $(0,\infty)$ with
\begin{equation}\label{eq:symcond}
f(x)=f(x/2)\quad\text{for almost every }x>0 .
\end{equation}
Consequently the naive ``$\mathbb E[Y\mid X=x]=\tfrac54x$ for every $x$'' is
incompatible with any proper prior.
\end{lemma}

\begin{proof}
Suppose $f$ satisfies \eqref{eq:symcond}, i.e.\ $f(t)=f(2t)$ for a.e.\ $t>0$;
iterating, $f(t)=f(2^{n}t)$ for all integers $n$ and a.e.\ $t$. Fix $a>0$ and let
$c:=\int_a^{2a} f\ge 0$. For each integer $n$, the substitution $x=2^{n}u$
(so $dx=2^{n}\,du$, and $f(x)=f(2^{n}u)=f(u)$) gives
\[
\int_{2^{n} a}^{2^{n+1}a} f(x)\,dx
=\int_{a}^{2a} f(2^{n}u)\,2^{n}\,du
=2^{n}\int_{a}^{2a} f(u)\,du
=2^{n}c .
\]
Splitting $(0,\infty)$ into the disjoint dyadic blocks $(2^{n}a,2^{n+1}a]$,
$n\in\mathbb Z$, gives
\[
\int_0^\infty f=\sum_{n\in\mathbb Z} 2^{n}c .
\]
If $c>0$ this series diverges (already its $n\ge 0$ tail does); if $c=0$ for all
$a$ then $f=0$ a.e. In neither case is $\int_0^\infty f=1$. Hence no probability
density satisfies \eqref{eq:symcond}.
\end{proof}

\begin{lemma}[Switching cannot be favorable for all $x$]\label{lem:notall}
For any probability density $f$ with $\mathbb E[T]<\infty$ it is impossible that
$2f(x)>f(x/2)$ for almost all $x>0$; equivalently, by Lemma~\ref{lem:naive},
switching cannot be strictly favorable at almost every observed value. Indeed,
for any such proper prior the expected gain from switching is exactly zero:
\begin{equation}\label{eq:meanzero}
\mathbb E[Y-X]=0 .
\end{equation}
\end{lemma}

\begin{proof}
Both envelopes contain the random amounts $T$ and $2T$, and the agent holds each
side with probability $\tfrac12$ independently of $T$. Hence, when
$\mathbb E[T]<\infty$ so that $X,Y$ are integrable,
\[
\mathbb E[X]=\tfrac12\,\mathbb E[T]+\tfrac12\,\mathbb E[2T]=\tfrac32\,\mathbb E[T],
\qquad
\mathbb E[Y]=\tfrac12\,\mathbb E[2T]+\tfrac12\,\mathbb E[T]=\tfrac32\,\mathbb E[T],
\]
so $\mathbb E[X]=\mathbb E[Y]$ and \eqref{eq:meanzero} holds. The first claim
follows: if $2f(x)>f(x/2)$ held for a.e.\ $x$, then by Lemma~\ref{lem:naive}
$\mathbb E[Y\mid X=x]>x$ for a.e.\ $x$, and integrating against the law of $X$
(which is licit because $X,Y$ are integrable, so the law of total expectation
$\mathbb E[Y]=\mathbb E[\mathbb E[Y\mid X]]$ applies) would force
$\mathbb E[Y]>\mathbb E[X]$, contradicting \eqref{eq:meanzero}.
\end{proof}

\begin{remark}\label{rem:proper}
Switching \emph{can} be advantageous for some observed values and disadvantageous
for others; what is impossible under a proper integrable prior is uniform
advantage. By Lemma~\ref{lem:naive} switching helps exactly where
$2f(x)>f(x/2)$ and hurts where $2f(x)<f(x/2)$; the two effects cancel in
expectation by \eqref{eq:meanzero}. (For $T$ exponentially distributed, $f(t)=e^{-t}$, the sign of
$2f(x)-f(x/2)=2e^{-x}-e^{-x/2}$ changes at $x=2\ln 2$: switching is favorable
for $x<2\ln2$ and unfavorable for $x>2\ln2$, with global mean gain $0$.) There is therefore no paradox: the optimal action
depends on $x$, and ``always switch'' is never optimal.
\end{remark}

\subsection{The improper-prior version}

\begin{lemma}[Conditional symmetry forces an improper prior with undefined
expectation]\label{lem:improper}
For a prior given by a density $f$ on $(0,\infty)$, the conditional symmetry
\eqref{eq:posterior}$=(\tfrac12,\tfrac12)$ for every $x$ holds if and only if
\begin{equation}\label{eq:fdyad}
f(x)=f(x/2)\quad\text{for (a.e.)\ }x>0 .
\end{equation}
\emph{Every} $\sigma$-finite measure with such a density is improper
(non-normalizable), and under it the relevant expectations diverge,
$\mathbb E[X]=\mathbb E[Y]=+\infty$, so that $\mathbb E[Y-X]=\infty-\infty$ is
undefined. The canonical such prior is the scale-invariant (log-uniform)
measure $d\mu(t)=c\,dt/t$ \((c>0)\), equivalently $T=e^{S}$ with $S$
translation-invariant on $\mathbb R$; but it is only one representative, the
condition \eqref{eq:fdyad} being satisfied by \emph{any} multiplicatively
$2$-periodic density.
\end{lemma}

\begin{proof}
\emph{Characterization \eqref{eq:fdyad}.} By \eqref{eq:posterior} the two
posterior probabilities are $f(x)/(f(x)+f(x/2))$ and $f(x/2)/(f(x)+f(x/2))$;
these equal $(\tfrac12,\tfrac12)$ precisely when $f(x)=f(x/2)$. (In the
two-hypothesis Bayes rule the parent pairs $(x,2x)$ and $(x/2,x)$ carry weights
proportional to $f(x)$ and $f(x/2)$.) This gives \eqref{eq:fdyad}.

\emph{Every such prior is improper.} Condition \eqref{eq:fdyad} is exactly
\eqref{eq:symcond}, so by Lemma~\ref{lem:noproper} no density satisfying it is
normalizable: writing $c:=\int_a^{2a}f\ge 0$ for a fixed $a>0$, the dyadic-block
computation in the proof of Lemma~\ref{lem:noproper} gives
$\int_0^\infty f=\sum_{n\in\mathbb Z}2^{n}c$, which is $+\infty$ unless $f\equiv0$.
We emphasize that \eqref{eq:fdyad} does \emph{not} single out the scale-invariant
prior: it holds for \emph{every} multiplicatively $2$-periodic $f$, i.e.\ every
$f$ of the form $f(t)=\varphi(\log_2 t\bmod 1)$ with $\varphi\ge 0$. The
log-uniform $d\mu(t)=c\,dt/t$ (for which $f(t)=c/t$, a particular function of
$\log_2 t$) is merely the canonical example. Thus no \emph{uniqueness} is
claimed or needed; what matters is that the entire family is improper.

\emph{Divergence of the expectations.} Let $f\ge 0$ satisfy \eqref{eq:fdyad} and
be not a.e.\ $0$, so $C:=\int_a^{2a} t\,f(t)\,dt>0$ for some $a>0$. The induced
measure of the observed value $X$ is, by \eqref{eq:Xlaw} and the two-hypothesis
decomposition, $g(x)\,dx$ with $g(x)=\tfrac12\bigl(f(x)+f(x/2)\bigr)\ge\tfrac12 f(x)$.
Hence, using the substitution $x=2^{n}u$ (so $f(x)=f(u)$, $dx=2^{n}du$, and
$x=2^{n}u$) on each dyadic block,
\[
\mathbb E[X]=\int_0^\infty x\,g(x)\,dx\ \ge\ \tfrac12\int_0^\infty x\,f(x)\,dx
=\tfrac12\sum_{n\in\mathbb Z}\int_{2^{n}a}^{2^{n+1}a} x\,f(x)\,dx
=\tfrac12\sum_{n\in\mathbb Z} 4^{n}\!\int_a^{2a}\! u\,f(u)\,du
=\tfrac12\sum_{n\in\mathbb Z}4^{n}C=+\infty .
\]
By the identical computation applied to the other envelope (the role of $f(x)$
and $f(x/2)$ in $g$ is symmetric and $Y$ ranges over the same set of values),
$\mathbb E[Y]=+\infty$ as well. Therefore $\mathbb E[Y]-\mathbb E[X]=\infty-\infty$
is undefined.

For the canonical case $d\mu(t)=c\,dt/t$ one has $g(x)=\tfrac12(c/x+2c/x)=\tfrac{3c}{2x}$
and the conditional probabilities \eqref{eq:posterior} equal $(\tfrac12,\tfrac12)$
for every $x$, so \eqref{eq:condexp} gives $\mathbb E[Y\mid X=x]=\tfrac54x>x$
everywhere, while $\mathbb E[X]=\int_0^\infty x\cdot\tfrac{3c}{2x}\,dx=+\infty$,
illustrating the general statement.
\end{proof}

\begin{proposition}[Resolution of the envelope paradox]\label{prop:env}
The ``always switch / always switch back'' contradiction never arises for a
coherent Bayesian.
\begin{enumerate}
\item For a \emph{proper} integrable prior, \eqref{eq:posterior}--\eqref{eq:condexp}
give the genuine posterior; by Lemma~\ref{lem:naive} the naive $\tfrac54x$ holds
only where $f(x)=f(x/2)$, by Lemma~\ref{lem:noproper} this cannot hold for all
$x$, and by Lemma~\ref{lem:notall} the expected gain from switching is $0$.
Switching is optimal only at those $x$ where $2f(x)>f(x/2)$, never universally;
the symmetry argument for ``switch back'' is then simply false because the agent
now conditions on a different observed value.
\item For the \emph{improper} scale-invariant prior, conditional symmetry does
hold, but by Lemma~\ref{lem:improper} the expectations are infinite and
$\mathbb E[Y-X]$ is undefined. The step ``$\mathbb E[Y\mid X=x]>x$ for all $x$,
therefore $\mathbb E[Y]>\mathbb E[X]$, therefore switch'' is invalid: the law of
total expectation $\mathbb E[Y]=\mathbb E[\mathbb E[Y\mid X]]$ requires
integrability (Fubini--Tonelli), which fails. Pointwise positivity of a
conditional expectation does \emph{not} imply a positive---or even defined---
unconditional expectation when the integrand is non-integrable.
\end{enumerate}
In both cases the decision-theoretic mandate---maximize a well-defined expected
utility---either makes ``always switch'' false (case 1) or makes the whole
computation inadmissible (case 2).
\end{proposition}

\begin{proof}
Immediate from Lemmas~\ref{lem:naive}, \ref{lem:noproper}, \ref{lem:notall} and
\ref{lem:improper}, together with the failure of the law of total expectation
for non-integrable variables: if $Z\ge 0$ has $\mathbb E[Z]=\infty$ one cannot
deduce a strict inequality between two infinite means. The contradiction in the
naive argument is precisely the attempt to apply, to the case-2 improper prior,
the sign-preserving identity $\mathbb E[Y]=\mathbb E[\mathbb E[Y\mid X]]$
that is only valid in case~1, where, however, the hypothesis
$\mathbb E[Y\mid X=x]>x\ \forall x$ is false.
\end{proof}

\section{The necktie paradox}

\subsection{Setup} Let the two tie prices be $A$ (mine) and $B$ (the other
person's), modeled as random variables with a joint prior $\pi$ on
$(0,\infty)^2$, with $\Pr_\pi(A=B)=0$ (ties of exactly equal price are a null
event, harmless to exclude). The agreed wager: the owner of the dearer tie
surrenders it to the other; equivalently each player's net monetary change is
\begin{equation}\label{eq:tiepayoff}
G_A=\begin{cases} +B & \text{if }A<B\quad(\text{I receive the dearer tie } B),\\
-A & \text{if }A>B\quad(\text{I surrender my dearer tie } A),\end{cases}
\qquad G_B=-G_A .
\end{equation}
(The transfer is pure: $G_A+G_B=0$ pointwise.)

\subsection{The naive reasoning and its defect} Player~$A$ observing
$A=v$ argues: ``with probability $\tfrac12$ mine is dearer and I lose $v$; with
probability $\tfrac12$ mine is cheaper and I gain $B>v$; hence
$\mathbb E[G_A\mid A=v]\ge \tfrac12(-v)+\tfrac12\,(\text{something}>v)>0$.'' The
defect is, again, the conditional symmetry assumption
$\Pr(A>B\mid A=v)=\Pr(A<B\mid A=v)=\tfrac12$ for \emph{every} $v$, together with
the assumption that, conditional on $A<B$, the conditional mean of $B$ exceeds
$v$. We show this cannot hold for all $v$ under any proper prior, and that the
correct, well-defined object is $\mathbb E_\pi[G_A]$, which is the negative of
$\mathbb E_\pi[G_B]$ and hence cannot be positive for both players.

\begin{theorem}[Zero-sum identity / impossibility of mutual advantage]\label{thm:fubini}
Let $\pi$ be any prior on $(0,\infty)^2$ under which $A$ and $B$ are integrable
(so $\mathbb E_\pi A<\infty,\ \mathbb E_\pi B<\infty$). Then $G_A$ and $G_B$
are integrable, $\mathbb E_\pi[G_A]+\mathbb E_\pi[G_B]=0$, and therefore it is
impossible that $\mathbb E_\pi[G_A]>0$ and $\mathbb E_\pi[G_B]>0$ simultaneously.
Moreover, if $\pi$ is exchangeable in the two coordinates
(\,$\pi(A\in S,B\in U)=\pi(A\in U,B\in S)$ for all measurable $S,U$\,), then
\begin{equation}\label{eq:tiezero}
\mathbb E_\pi[G_A]=\mathbb E_\pi[G_B]=0 ,
\end{equation}
so the wager is fair for each player.
\end{theorem}

\begin{proof}
Since $|G_A|\le\max(A,B)\le A+B$ and $A,B$ are integrable, $G_A$ (and
$G_B=-G_A$) is integrable; in particular $\mathbb E_\pi[G_A]$ is a finite,
well-defined number. From \eqref{eq:tiepayoff}, $G_A+G_B=0$ pointwise (if $A>B$
then $G_A=-A$, $G_B=+A$; if $A<B$ then $G_A=+B$, $G_B=-B$), so linearity of the
well-defined expectation gives
$\mathbb E_\pi[G_A]+\mathbb E_\pi[G_B]=\mathbb E_\pi[G_A+G_B]=0$. Two real numbers
summing to zero cannot both be strictly positive, proving the impossibility.

For the exchangeable case, the swap map $\sigma:(a,b)\mapsto(b,a)$ pushes $\pi$
to itself and, by \eqref{eq:tiepayoff}, sends $G_A$ to $G_B$ (it exchanges the
roles of the two players); hence
$\mathbb E_\pi[G_A]=\mathbb E_\pi[G_A\circ\sigma]=\mathbb E_\pi[G_B]$. Combined
with $\mathbb E_\pi[G_A]+\mathbb E_\pi[G_B]=0$ this yields \eqref{eq:tiezero}.
\end{proof}

\begin{proposition}[Resolution of the necktie paradox]\label{prop:tie}
The conclusion ``both players find the wager strictly favorable'' is impossible
for coherent Bayesians.
\begin{enumerate}
\item If both players use proper priors with integrable prices, then by
Theorem~\ref{thm:fubini} their well-defined expected gains sum to zero, so they
cannot both be positive. If, in addition, beliefs are exchangeable (the natural
symmetric situation, where neither has private information distinguishing the
ties), each expected gain is exactly $0$ and the bet is fair, not favorable.
\item The naive argument obtains a positive expectation for both only by
asserting, for every observed $v$, both the conditional symmetry
$\Pr(A>B\mid A=v)=\tfrac12$ \emph{and} that the conditional mean of the
opponent's price on the event $\{B>v\}$ exceeds $v$ by more than the symmetric
loss term, so that $\mathbb E[G_A\mid A=v]>0$ for every $v$. As shown in the
proof of part~(1) below, $\mathbb E[G_A\mid A=v]>0$ for a.e.\ $v$ is impossible
under any integrable exchangeable prior, since it would force
$\mathbb E_\pi[G_A]>0$ in contradiction with \eqref{eq:tiezero}. Retaining it for
every $v$ therefore requires an improper prior under which
$\mathbb E_\pi[\max(A,B)]=\infty$ and $\mathbb E_\pi[G_A]$ is the undefined form
$\infty-\infty$ --- the two-dimensional analogue of Lemma~\ref{lem:improper}. In
that regime the law of total expectation again
fails and ``$\mathbb E[G_A\mid A=v]>0\ \forall v \Rightarrow \mathbb E[G_A]>0$''
is invalid.
\end{enumerate}
\end{proposition}

\begin{proof}
Part (1) is Theorem~\ref{thm:fubini}. For part (2): conditionally,
\[
\mathbb E[G_A\mid A=v]
=\mathbb E[B\,\mathbf 1_{B>v}\mid A=v]-v\,\Pr(B<v\mid A=v).
\]
The naive evaluation replaces $\Pr(B<v\mid A=v)$ by $\tfrac12$ and asserts the
first term exceeds $\tfrac12 v$ for every $v$. Suppose, for contradiction, that
$\mathbb E[G_A\mid A=v]>0$ for a.e.\ $v$ under an exchangeable prior with
integrable prices. Then integrating against the marginal of $A$---legitimate
because $A,B$ are integrable, so the law of total expectation
$\mathbb E[G_A]=\mathbb E[\mathbb E[G_A\mid A]]$ holds---gives
$\mathbb E_\pi[G_A]>0$, contradicting \eqref{eq:tiezero}. Hence such uniform
conditional positivity is impossible under any integrable exchangeable prior:
\emph{this is the decisive point}, and it does not depend on any special form of
the prior. The only way to retain ``$\mathbb E[G_A\mid A=v]>0$ for every $v$''
is therefore to abandon integrability and use an improper prior under which
$\mathbb E_\pi[\max(A,B)]=\infty$ (e.g.\ a scale-invariant joint prior of the
log-uniform type, the two-dimensional analogue of Lemma~\ref{lem:improper}).
Then $\mathbb E_\pi[G_A]$ is the undefined form $\infty-\infty$, and the
inference ``conditional gain $>0$ everywhere $\Rightarrow$ favorable'' fails
because $\mathbb E[G_A]=\mathbb E[\mathbb E[G_A\mid A]]$ requires the integrability that
now does not hold.
\end{proof}

\section{The unified principle}

\begin{theorem}[Unified resolution]\label{thm:unified}
Both paradoxes are instances of the same error and are resolved by the same
principle.
\begin{enumerate}
\item \textbf{Located error.} The fallacious step is the assignment, by the
principle of indifference, of probability $\tfrac12$ to each of the two states
(``mine is larger'' / ``mine is smaller'') \emph{conditional on the observed
value, for every observed value}. The principle of indifference governs prior
symmetry, not posterior probabilities; the genuine posterior \eqref{eq:posterior}
depends on the prior density and equals $(\tfrac12,\tfrac12)$ for all
observations only under an improper scale-invariant prior.
\item \textbf{Why it fails.} (i) For a proper integrable prior the conditional
symmetry is false (Lemmas~\ref{lem:naive}, \ref{lem:noproper}), and the correct
expected gain from exchange is $0$ (Lemma~\ref{lem:notall},
Theorem~\ref{thm:fubini}). (ii) For the improper prior that does make the
conditional symmetry hold, the expected payoff is non-integrable:
$\mathbb E[\,\cdot\,]=\infty-\infty$ is undefined and the law of total
expectation $\mathbb E[\,\cdot\,]=\mathbb E[\mathbb E[\,\cdot\mid\text{obs}]]$
fails (Lemma~\ref{lem:improper}). Thus the chain ``conditional gain $>0$
everywhere $\Rightarrow$ unconditional gain $>0$ $\Rightarrow$ act'' is never
simultaneously sound and applicable.
\item \textbf{Normative principle.} A coherent subjectivist must base the
exchange decision on $\mathbb E_\pi[\,U(\text{after})-U(\text{before})\,]$ where
(a) $\pi$ is a prior under which this quantity is well defined and finite
(equivalently, the relevant monetary variables are $\pi$-integrable, so that
Fubini--Tonelli legitimizes conditioning), and (b) all conditional probabilities
are the true posteriors $\pi(\cdot\mid\text{observed data})$, never values
imposed by symmetry. Under this principle, whenever $\pi$ is symmetric under the
interchange of the two parties, the expected change in utility from the symmetric
exchange is zero (Lemma~\ref{lem:notall}, Theorem~\ref{thm:fubini}); hence no
symmetric exchange can be advantageous to both parties, and ``always switch'' is
never optimal for all observations.
\end{enumerate}
\end{theorem}

\begin{proof}
Parts (1)--(2) collect Propositions~\ref{prop:env} and \ref{prop:tie}. For (3):
when $\pi$ is exchangeable and the monetary variables are integrable, the
exchange map $\sigma$ (swap envelopes, resp.\ swap ties) preserves $\pi$ and
negates the net transfer, so
$\mathbb E_\pi[\Delta U]=\mathbb E_\pi[\Delta U\circ\sigma]=-\mathbb E_\pi[\Delta U]$,
whence $\mathbb E_\pi[\Delta U]=0$; this is the content of \eqref{eq:meanzero}
and \eqref{eq:tiezero}. Integrability is exactly the condition that makes this
expectation, and the law of total expectation used to pass between conditional
and unconditional statements, valid---precisely the hypothesis that fails in the
improper-prior versions. Therefore adherence to the stated principle excludes the
paradoxical conclusions in every exchange-type problem of this form.
\end{proof}

\begin{remark}[Role of nonlinearity of utility]\label{rem:utility}
We assumed linear utility to isolate the probabilistic error. If utility $u$ is
strictly concave (risk aversion) the same analysis applies to $u$ in place of the
identity: the exchange decision is governed by $\mathbb E_\pi[u(Y)-u(X)]$, which
under an exchangeable proper prior is still $0$ by the symmetry argument of
Theorem~\ref{thm:unified}(3) (the swap map negates $u(Y)-u(X)$). Thus
nonlinearity of utility is \emph{not} the source of the paradox---the resolution
is the same with or without it; the source is the illegitimate
conditional-symmetry assumption and, in the limiting case, the use of a
non-integrable improper prior.
\end{remark}

\end{document}
